\section{Implementation}
\label{sec:implementation}

The prototype maps the design onto a modified Linux kernel with a
\code{cgroup/namei_ext} attachment point. Policies are eBPF programs under
\file{bpf/policies/*.bpf.c}. The implementation consists of the VFS call
sites, the BPF ABI, user-space policy attachment and map setup, and kernel-side
validation for the supported action set.

\subsection{Kernel Attachment}

The lookup path calls the policy before the VFS commits to the lower object for
the current component or final-component create/open path. Directory
enumeration calls the same decision function for entries that may be exposed
under an alias. Both call sites preserve the normal lower-filesystem path after
the policy returns a validated decision, so the VFS still applies ordinary
permission checks and file operations to the selected object.

\subsection{BPF ABI}

The kernel-facing ABI exposes one decision function. The context includes an
event type, operation flags, the current component name, parent identity,
policy-relevant identifiers, and bounded output fields for the selected
path-view action. The current action set is \code{PASS}, \code{REDIRECT},
\code{HIDE}, and \code{SELECT\_TARGET}. \code{REDIRECT} carries a bounded
replacement component in \code{redirect_name}. \code{SELECT\_TARGET} carries an
opaque \code{target_id} that indexes a kernel-held registered \code{struct
path}. The policy cannot return arbitrary unbounded path strings or perform
recursive path walks.

The kernel validates action type, bounded names, target identity, and selection
scope before continuing. Redirect decisions are validated for length and
same-parent scope. Hide returns absence on lookup and suppresses the lower
entry during directory enumeration.

The current prototype supports registered target selection only for directory
targets and does not implement final-file target selection.

\subsection{Prototype Coverage}

The prototype implements pass-through, validated same-parent
replacement-component redirect, hide, and an initial registered directory-target
selection path. The registered directory-target path supports intermediate directory
components and final directory operations such as stat, \code{O\_DIRECTORY}
open, and readdir over the selected lower directory.

The prototype fails closed for create, non-directory final opens, final-file
target selection, and synthetic parent-directory aliases.

\subsection{Replacement-Component Validation And Permission Preservation}

Workload setup configures the policy state that chooses replacement components.
The prototype implements the redirect subset by copying a bounded
\code{redirect_name} from the BPF output, validating its length, action, and
same-parent scope, and then resuming the normal VFS path. The replacement
component is not an
arbitrary path string and does not trigger a recursive path walk. After
validation, the normal VFS path continues, including permission checks and
lower-filesystem file operations. Same-parent validation keeps \namei from
becoming arbitrary path rewriting or a filesystem service.

The implemented hide action uses the same decision point but does not allocate
or synthesize filesystem objects. Lookup observes absence, and directory
enumeration suppresses the lower entry for the active policy state.

Registered directory-target selection adds kernel-owned object identity and lifetime
validation before a policy can select a lower directory outside the visible
parent.

\subsection{Failure Handling}

Unsupported actions, verifier failures, malformed decisions, syscall failures,
and invalid lower selections are hard failures in the owning Make target. The
implementation does not silently fall back to a weaker policy when validation
fails. This fail-closed behavior is part of the safety boundary evaluated in
RQ3.

Phase 1 validation runs in KVM with the modified kernel and the real
\code{cgroup/namei_ext} attach path. Host-only checks, object inspection, and
bpftool-only smoke tests are diagnostics rather than validation for the paper's
mechanism claims.
